\section{Threat model}
\label{sec:threat-model}

This section describes how keystroke logs of writing are released (\S\ref{sec:threat-model:releases}), the kind of
auxiliary information that this study considers (\S\ref{sec:threat-model:auxiliary}) and who holds it
(\S\ref{sec:threat-model:settings}). It then states the attacker (\S\ref{sec:threat-model:attacker}) and the direction
in which the assumptions bias the results (\S\ref{sec:threat-model:bias}). The attacker's knowledge, the candidate sets
and the list of attackers were fixed in the preregistration of this study before any data were acquired; this section
motivates them and does not change them.

\subsection{Releases of keystroke logs}
\label{sec:threat-model:releases}

\paragraph{Logs are released with the text hidden.}
Keystroke logs of writing are released to support research on the writing process, and a release may hide the text
and keep the rest of the log. The Kaggle competition ``Linking Writing Processes to Writing Quality'' released logs of
argumentative essays in which, in the words of its data description, ``to prevent reproduction of the essay text, all
alphanumeric character inputs have been replaced with the `anonymous' character q;\ punctuation and other special
characters have not been anonymized''~\cite{kaggle2023lwp}. The released columns include the press and release times,
the activity, the text change, the cursor position and the word count. This is the release that motivates the study:
the text is hidden, and every letter and digit is still present as a press.

\paragraph{The lengths that remain are used.}
Masking each letter and digit keeps the length of every word. The winning solution of the competition began by
rebuilding the essay from the masked log (``the first thing I did in this competition is to reconstruct sentences from
keyboard activities'') and used statistics of the rebuilt text, such as words per sentence and word length, to predict
the score~\cite{inubushi2024first}. The structure that makes a masked log useful for research on writing is also the
structure that a party holding candidate texts can compare with those texts.

\paragraph{Other masking keeps lengths too.}
Two public datasets of keystrokes in an introductory programming course mask the characters of identifying text, such
as a student's name, with the character \texttt{@}; a~function \texttt{even()} in which the name ``Eve'' was masked
becomes \texttt{@@@n()}~\cite{edwards2023csedm}. This release hides names, not the program, so it is not an instance of
the threat studied here; it shows that masking one character for one character is a common choice, and with it the
lengths of the masked strings remain. The times of the events are not processed in this release.

\paragraph{Parts of logs.}
A log may also be shared or analysed only in part. The evaluation therefore covers windows of a log as well as full
logs (Section~\ref{sec:protocol}).

\subsection{Auxiliary information other than labelled keystrokes}
\label{sec:threat-model:auxiliary}

Earlier work on releasing keystroke data considered an attacker who holds a labelled sample of the same person's
keystrokes. Keystroke latencies link students across samples, and rounding the latencies in typing profiles reduced
identification~\cite{leinonen2017preventing}. The paper that released the programming datasets answers the risk of
identification by timing in the same terms: ``all of these identification techniques require a labeled sample, which
the regulatory definition of identifiability assumes is not available'', and ``the data is deidentified only as long
as there is no other labeled sample of a student's keystrokes''~\cite{edwards2023csedm}. A survey of anonymisation for
behavioural biometric data notes that ``most anonymization techniques are only evaluated rudimentarily with the
assumption of a weak attacker''~\cite{hanisch2021survey}.

This study considers a different kind of auxiliary information: the texts. Records in a release can be linked to
outside information that the releasing party did not anticipate~\cite{narayanan2008robust}, and for logs of writing the
most direct such information is the text that was written. A log matched to a text whose writer is known would, from
then on, also be a labelled sample of that writer's keystrokes.

\subsection{Who holds the texts}
\label{sec:threat-model:settings}

A hidden text can be matched only if it exists somewhere outside the release. Table~\ref{tab:settings} lists the
settings we consider. They differ in what the holder has, in the direction of the search and in whether the true text
is always among the candidates (the closed world assumed by the evaluation).

\begin{table*}
  \caption{Settings in which texts are held outside a release of processed logs.}
  \label{tab:settings}
  \small
  \begin{tabular}{>{\raggedright\arraybackslash}p{0.15\linewidth}>{\raggedright\arraybackslash}p{0.19\linewidth}>{\raggedright\arraybackslash}p{0.12\linewidth}>{\raggedright\arraybackslash}p{0.19\linewidth}>{\raggedright\arraybackslash}p{0.21\linewidth}}
    \toprule
    Setting & Holder has & Direction & Closed world holds when & Newly learned by matching \\
    \midrule
    Public corpus on the same prompts & Texts of a public corpus collected in the same way & One log, many texts &
      The hidden essays are in the corpus & The hidden text, and the score and process columns of its log \\
    Classroom & The texts and names of the students of a class or course & One log, many texts & The release consists
      of that class & The writing process, given for research, of a named student \\
    Published essay & One text and its writer, found in a class or online & One text, many logs & Rarely: the log need
      not be in the release (open world) & The writing process and score of a known writer \\
    \bottomrule
  \end{tabular}
\end{table*}

\paragraph{A public corpus on the same prompts.}
The KLiCKe corpus was collected on Amazon Mechanical Turk in 2022 and released together with its keystroke logs and
human scores: 4,992 argumentative essays on four prompts derived from the SAT, each written in 30
minutes~\cite{tian2025klicke}. Because its texts are public, its logs need no protection of the text; it is the
data on which we evaluate the processing. The description of the procedure of the competition data names the same
platform, the same kind of prompts and the same sequence of tasks~\cite{kaggle2023lwp}. A public set of about a
thousand essays per prompt is therefore the kind of candidate set that a text-hidden release on the same prompts would
face. The closed world holds only if the hidden essays are among the public ones. We did not examine whether the two
releases share essays, for ethical reasons (see the ethical considerations).

\paragraph{Classrooms.}
Teachers and course staff hold the texts submitted by their students, with names. The programming datasets above were
collected in a course~\cite{edwards2023csedm}, and writing tools now keep the text and the writing process of students
in one product and offer a replay of the process to instructors~\cite{grammarly2026authorship,turnitin2026clarity}.
We know of no release of logs from such tools for research. If a release contained the logs of one class with the
text processed, its teacher would hold every candidate; if the release spans many classes, the teacher holds the
candidates for only part of it, and the closed world does not hold.

\paragraph{A published essay.}
An essay may be published by its writer, in a class or online. A party who finds it holds one text and searches a
release for its log. The search runs in the other direction, and the log may not be in the release at all. The
evaluation does not run this direction, and its closed-world accuracy is only a guide in this setting: with many logs
to search and a true log that may be absent, false matches are expected to be more frequent.

\paragraph{Translation.}
In research on the translation process, translations of the same source text by other translators are held by the
researchers who collected them. The Japanese translation logs evaluated in this study use these translations as
candidates (Section~\ref{sec:results}).

\subsection{Attacker}
\label{sec:threat-model:attacker}

\paragraph{Knowledge.}
The attacker holds one processed log and $N$ candidate texts that include the true text, all written on the same
prompt. The attacker knows the processing procedure, following Kerckhoffs's principle~\cite{kerckhoffs1883}. The
attacker may train on public data of the same kind, that is, logs before processing paired with their texts; in this
study these are the writers of the prompts that are not being evaluated, and their texts are never candidates. Beyond
the processed log and the candidates, the attacker does not use the writer's earlier texts, text generated by language
models, or links across several releases (such as several logs of the same writer). In the taxonomy of Wagner and
Eckhoff~\cite{wagner2018metrics}, the attacker is passive, external (it observes the release from outside the
releasing party) and static, and has prior knowledge (the candidate set and the training data). One attacker in this
study, the public KeyTAR model~\cite{qiu2026keytar}, contains a public language model; it is reported as a reference
outside this threat model.

\paragraph{Goals and harms.}
The main harm is (b): linking a log and its other columns (times, pauses, the revision process, pasting, the score)
to a text and, where the writer of the text is known, to the writer. In the classroom and published-essay settings,
this ties a record that a writer gave for research to a named person. A secondary harm is (a): learning which text a
processed log hides, which defeats the purpose of the processing even when that text is public elsewhere. Both rest on
the link between a log and its text, which is what this study measures; the study does not identify writers. Because
the protection given by a processing step may differ between writers, the distribution of accuracy per writer is
reported as well as the mean.

\paragraph{Size of the candidate set.}
We use $N = 10$, $100$ and $1{,}000$. The value $1{,}000$ is the size of a public corpus on one prompt; $10$ and
$100$ are illustrative sizes of a class or a list of submissions and of a year or a course. $N$ gives only the number
of candidates. The decoys are drawn uniformly from writers on the same prompt, who need not be as similar to one another
as the members of one class.

\subsection{Assumptions and the direction of bias}
\label{sec:threat-model:bias}

Some assumptions favour the attacker: the true text is always among the candidates (the closed world, examined
otherwise only as a diagnostic), the attacker trains on data collected in the same way as the evaluated logs, and in
windows the attacker knows where the window starts and ends. Others limit it: the attackers do not use language models,
the writer's earlier texts or several releases, because the study measures what the processed columns themselves leave.
The results therefore describe the attackers tried here, under these assumptions. They give no proof of safety and no
upper bound on the risk of matching; a processing step that defeats these attackers may not defeat a stronger one.

Under this threat model a releasing party chooses which columns to keep, how coarsely to give times and lengths, and
whether to release full logs or parts of them. The results are material for that choice; they do not recommend a
form.
